\documentclass{article}

% NeurIPS style file
\usepackage[final]{neurips_2024}

% Packages
\usepackage[utf8]{inputenc}
\usepackage[T1]{fontenc}
\usepackage{hyperref}
\usepackage{url}
\usepackage{booktabs}
\usepackage{amsfonts}
\usepackage{amsmath}
\usepackage{nicefrac}
\usepackage{microtype}
\usepackage{xcolor}
\usepackage{graphicx}
\usepackage{placeins}

\usepackage{tabularx}
\usepackage{array}
\usepackage{float}

\title{ Quantitative Trading Optimization: A Machine Learning Approach of Stock Market Forecasting and Trading Decision Making via Hybrid LSTM and DQN}

\author{
Langning Li \& Runpeng Zhu \\
University of Rochester \\
Rochester, NY 14627, USA \\
\texttt{lli81@u.rochester.edu \quad rzhu12@u.rochester.edu }
}

\begin{document}

\maketitle

\begin{center}
\fbox{\parbox{0.92\linewidth}{\small\textbf{Reproducibility notice (September 2026).} The numerical tables and output cells in this historical report were generated before the repository adopted the canonical leakage-aware evaluation contract. The ARIMA and vanilla-LSTM baselines used a one-day target while the hybrid model used a 20-day target; the hybrid threshold was also fit before the chronological split. The DQN evaluation normalized signals with evaluation-period statistics and omitted trading frictions. These values are retained for provenance only and must not be presented as current apples-to-apples evidence. See \texttt{README.md} and \texttt{METHODOLOGY.md} for the corrected protocol.}}
\end{center}

\begin{abstract}
Stock market forecasting is a challenging machine learning problem because financial data are noisy, nonlinear, and affected by external factors. In this project, we develop a two-stage framework that combines a multivariate LSTM forecaster with a Deep Q-Network (DQN) trading agent. We evaluate forecasting with directional and probability-based metrics and evaluate trading with risk-adjusted, cost-aware portfolio metrics. The repository now defines a shared leakage-aware target, chronological split, train-fitted transformations, and exposure-limited backtest; the historical model outputs are retained for provenance and must be regenerated under this contract before performance conclusions are drawn.

\end{abstract}

\section{Introduction}
Financial market is difficult to predict because the change of the price is being affected by many interacting factors [4]. Historical prices and returns sequence contain useful temporal information, but they are often noisy and unstable. Also, the broader external factors such as macroeconomic conditions, interest rates, and geopolitical risk can also influence how the market price might change. This makes stock forecasting a challenging task to do, especially when the final objective have to both predict the future values of the stock and making profitable trading decisions.

A simple method for financial forecasting is to use the traditional statistical time series models. These models are not very complicated and can be easily understand, but they have a hard time capturing nonlinear relationships in market data. Deep learning models, especially the recurrent neural networks such as LSTM, provide a better option because they can learn and remember the temporal patterns from sequential data. However, using only historical returns may ignore important external information that affects market movement. Therefore, combining market return sequences with additional financial and macroeconomic features may provide a more complete representation of the market state.

In this project, we try to address stock market forecasting and trading decision problem making by building a two stage hybrid framework. In the first stage, the forecasting model predicts future market movement using both historical return information and external features. In the second stage, the predicted signal is used by a trading agent to select actions in a simulated trading environment. By connecting prediction with decision making, this framework allows us to evaluate not only forecasting accuracy, but also how useful the predicted signals are for quantitative trading.

\section{Problem Formulation and Project Goal}

Let $d_t$ denote the market future market direction at time $t$. The price forecasting task is to predict the direction of the future market return using historical information. Given a sequence of past stock market data,
\[
X_t = \{x_{t-k+1}, x_{t-k+2}, \ldots, x_t\},
\]
where $k$ is the size of the sliding window, the model predicts the future stock market direction or the probability of whether the future market return will be positive.

In this project, the whole problem is being divided into two parts: market direction forecasting and trading decision making. The stage one's forecasting task is mainly being treated as a directional prediction problem. Instead of predicting the exact future stock price or return value, the canonical protocol predicts whether the 20-trading-day forward return exceeds a threshold fitted on training outcomes only. Using the historical return sequences and external financial features in the dataset, the forecasting model learns to estimate the future direction of market movement.

The second task is being treated as a strategy making problem. At each time step, the trading agent observes the market related features and the predicted market signal, then chooses a trading action. Under the canonical protocol, actions represent unlevered positions from $-100\%$ to $+100\%$, and turnover costs and slippage are charged explicitly. The final reward of the model is based on the resulting portfolio return, with an additional alignment component that encourages the agent to choose positions consistent with the predicted market direction. In this way, the project connects the stage one directional forecasting model and the stage two decision optimization model together.

\section{Methodology}

\subsection{Dataset}
The dataset used in this project is a multivariate NASDAQ time series dataset. It contains historical market features such as open, high, low, close price, trading volume, and daily return, as well as external financial and macroeconomic features  such as VIX, treasury yield, dollar index, federal funds rate, unemployment rate, and CPI-related variables. The canonical preparation in \texttt{quant\_pipeline} uses a 20-trading-day forward return, a chronological 70\%/10\%/20\% split over labelled anchors, and training-only target/scaler statistics. These features provide complete and useful information for both short term price dynamics and the market's long term trend.

\subsection{ARIMA Baseline}

\textit{Historical implementation note: this baseline was originally trained for one-step return forecasting and is not comparable with the corrected 20-day classification target until it is retrained.}

We use the ARIMA model as the first baseline for stage one: a traditional statistical time series model [1]. For this baseline, we only use the historical NASDAQ return sequence as input into the ARIMA, and we did not incorporate any of the other external financial features in the dataset into this part. The return is being computed from the NASDAQ everyday stock market price by percentage change, and the dataset being split by standard: first 80\% of the data used for training and the remaining 20\% used for testing.

In our implementation, we use an ARIMA (5,0,1) model with a walk forward forecasting setup. To make the experiment computationally manageable, we evaluate the model on the most recent 500 test points. At each test step, the model is fitted using the most recent 1000 historical return values and predicts the next return. After each prediction, the true return is added back into the history before moving to the next step. This simulates a realistic forecasting setting where only past information is available.

Finally, we convert the ARIMA prediction into a directional signal by checking whether the predicted return is positive or negative. A positive predicted return will be labeled as 1, while a negative predicted return will be labeled as 0. This allows the ARIMA baseline to produce a direction based output that is consistent with the stage one goal of predicting future market movement.

\begin{verbatim}
model = ARIMA(history[-1000:], order=(5, 0, 1))
model_fit = model.fit()
pred = model_fit.forecast(steps=1)
\end{verbatim}

\subsection{Vanilla LSTM Baseline}

\textit{Historical implementation note: this baseline was originally trained for one-step return forecasting and is not comparable with the corrected 20-day classification target until it is retrained.}

We use the vanilla LSTM model as the deep learning baseline for the stage one. Compared with ARIMA, the LSTM can learn nonlinear temporal patterns from sequential data [2]. However, this baseline still only uses the historical return sequence as input, and we did not incorporate any of the other external financial features in the dataset into this model's input.

For this model, we input the data into the model using a sliding window. Each sample contains the past 30 days of return values, and the model predicts the return at the next time step. The dataset being split by standard: the first 80\% of the data is used for training, the next 10\% is used for validation, and the remaining 10\% is used for testing.

The model contains two layers: a LSTM followed by a linear output layer. In our implementation, the hidden size is set to 256, dropout is set to 0.2, and the batch size is 32. The model uses the Adam optimizer with a learning rate of 0.001. Early stopping is applied based on validation loss, with a patience value of 20 epochs.

\begin{verbatim}
model = MyLSTMSequential(
    input_size=1,
    hidden_size=256,
    output_size=1,
    num_layers=2,
    dropout=0.2
)
optimizer = optim.Adam(model.parameters(), lr=0.001)
criterion = nn.MSELoss()
\end{verbatim}

\subsection{Hybrid LSTM Model}

The hybrid LSTM model is our self designed stock market value future forecasting model in this project. This hybrid model is based on the idea that LSTM models can learn useful temporal patterns from financial time series data [2, 6]. Unlike the ARIMA and vanilla LSTM baselines, this model uses a multivariate feature set instead of only the historical return sequence. The goal is to predict the probability of future market direction, or whether the market price is more likely to move up or down.

For the target label, we use the 20-day forward close-to-close return. Rather than predicting the exact future return value, we treat the task as a binary classification problem: whether that return exceeds the median forward return observed in the training split. The threshold is never computed from validation or test outcomes. The model outputs one logit, later converted into a probability of the threshold event.

The model uses a two-stream architecture. The first stream takes sequential market features with a 60-day sliding window, where each sample has shape $60 \times 30$. This sequence is processed by an LSTM with hidden a hidden units of size 32. The second stream takes a macro snapshot with 7 features from the most recent day of the window. These macro features are passed through a linear layer and ReLU activation to produce a 16-dimensional macro embedding. The two outputs are then concatenated and passed through a dropout layer and a final linear layer to produce the prediction logit.

\vspace{0.8 em}
\begin{verbatim}
Sequential features (60 days x 30)     Macro snapshot (7,)
             |                                  |
       LSTM (hidden=32)              Linear(7 -> 16) + ReLU
             |                                  |
    last hidden state (32,)          macro embedding (16,)
             \______________ concat ___________/
                            |
                Dropout(0.3) -> Linear(48 -> 1)
                            |
                          logit
\end{verbatim}
\vspace{0.8 em}

This design is motivated by the idea that daily market features and macroeconomic features behave differently [6]. Market returns and the other technical indicators can change quickly within a short amount of time, so they are being processed by a LSTM time series model. In contrast, macro and risk-regime variables usually change slowly, so they are passed into the model as a separate snapshot. So creating two separate pipelines that process the data differently preserves long term signal. Later, by combining these two separate streams, the model can utilize both short term price info and long term market signals for future stock market direction prediction.

The model is trained using \texttt{BCEWithLogitsLoss}, which combines a sigmoid operation with binary cross-entropy loss and is suitable for binary classification. We use the Adam optimizer with learning rate 0.001, dropout of 0.3, gradient clipping for training stability, and early stopping based on validation performance.


\begin{verbatim}
criterion = nn.BCEWithLogitsLoss()
optimizer = torch.optim.Adam(model.parameters(), lr=0.001)
scheduler = torch.optim.lr_scheduler.ReduceLROnPlateau(
    optimizer, mode="min", factor=0.5, patience=5
)

torch.nn.utils.clip_grad_norm_(model.parameters(), max_norm=1.0)
\end{verbatim}
   
\subsection{DQN-Based Trading Agent}

For the second stage, we use a Deep Q-Network (DQN) agent to convert the stage one hybrid LSTM prediction signal into trading strategies [3]. The goal of this part is to learn to use the predicted signal to make effective trading strategies under the trading environment. At each time step, the agent observes the LSTM probability signal, recent market information, current portfolio position, cash ratio, and market regime information, and used those observed data to choose the correct action from the discrete action space.

In the corrected implementation, the action space contains 21 possible actions. Each action represents a target portfolio position from $-100\%$ short exposure to $100\%$ long exposure. A negative position represents short selling; leverage is not used in the canonical comparison. Turnover cost and slippage are charged whenever the target position changes.

The reward function is mainly based on the portfolio return generated by the selected position. In addition, we add an alignment bonus that encourages the agent to choose positions consistent with the predicted market direction from the hybrid LSTM model. If the predicted signal suggests an upward movement, long positions are encouraged; if the signal suggests a downward movement, short or lower exposure positions are encouraged. The reward design is shown below:

\begin{verbatim}
portfolio_return = daily_return * pos_ratio
reward_base      = portfolio_return * 100

desired   = 2 * signal_t - 1
alignment = (1 + pos_ratio * desired) / 2
bonus     = alignment * abs(daily_return) * 20
trading_cost = turnover * (transaction_cost + slippage)

reward = reward_base + bonus - trading_cost
\end{verbatim}

The DQN model is trained with a length of 252 trading days, and this is approximately one trading year (exclude the weekends and holidays). We use a replay buffer, epsilon greedy exploration, and a target network update mechanism through this part. The policy network is a multilayer perceptron with two hidden layers of size 256. The model is trained for 1,000,000 timesteps with learning rate $10^{-4}$, discount factor $\gamma = 0.97$, batch size 128, and replay buffer size 100,000.

\begin{verbatim}
dqn = DQN(
    "MlpPolicy",
    env,
    buffer_size=100000,
    learning_starts=5000,
    batch_size=128,
    learning_rate=1e-4,
    gamma=0.97,
    target_update_interval=1000,
    policy_kwargs=dict(net_arch=[256, 256])
)
\end{verbatim}

\section{Experimental Results}

This section records the historical outputs for both stages of the project. The tables are explicitly marked as retired because the models used different targets, split rules, signal transformations, and trading assumptions. Fresh ARIMA, vanilla-LSTM, hybrid-LSTM, and DQN results should be added only after retraining and evaluating through the canonical contract.

\subsection{Stage - One Forecasting Results (Legacy)}

Table 1 shows the future market direction prediction of the three stage one models. Since our main goal is to predict whether the future market value will move up or down, we focus on classification-based metrics: accuracy, F1 score, and AUC-ROC.

\begin{table}[H]
\centering
\small
\caption{Historical stage-one results (retired; models used different targets and are not directly comparable).}
\label{tab:stage1_results}
\begin{tabular}{lccc}
\toprule
\textbf{Model} & \textbf{Accuracy} & \textbf{F1 Score} & \textbf{AUC-ROC} \\
\midrule
ARIMA & 0.5400 & 0.6417 & -- \\
Vanilla LSTM & 0.5531 & 0.7123 & -- \\
Hybrid LSTM & 0.6291 & 0.7069 & 0.6672 \\
\bottomrule
\end{tabular}
\end{table}

\subsection{Stage-Two Trading Results (Legacy)}

After evaluating the stage one forecasting models, we use the output of the stage one hybrid LSTM model as part of the input to the DQN trading agent. The goal of this part is to test whether the predicted market direction signal can be converted into useful trading decisions. The DQN agent is evaluated on the period that range from 2019 to 2026, with an initial portfolio value of \$10,000. We compare the DQN strategy with three basic strategy: buy \& hold strategy, always cash strategy, and a fixed threshold strategy.

\begin{table}[H]
\centering
\small
\caption{Historical stage-two trading results (retired; pre-cost and pre-contract).}
\label{tab:dqn_results}
\begin{tabular}{lcccc}
\toprule
\textbf{Strategy} & \textbf{Total Return} & \textbf{Sharpe} & \textbf{Max Drawdown} & \textbf{Win Rate} \\
\midrule
DQN & 510.90\% & 0.954 & 35.65\% & 51.14\% \\
Buy \& Hold & 182.67\% & 0.751 & 36.40\% & 55.65\% \\
Always Cash & 0.00\% & 0.000 & 0.00\% & 0.00\% \\
Fixed Threshold & 134.57\% & 0.699 & 34.97\% & 27.15\% \\
\bottomrule
\end{tabular}
\end{table}

\section{Analysis and Discussion}
\subsection{Stage - One Forecasting Results}
The numerical stage-one values in Table~\ref{tab:stage1_results} are retained as a historical record only. ARIMA and vanilla LSTM were trained for one-step return regression, while the hybrid model used a 20-day classification label; their accuracy and F1 values therefore do not form a valid comparison. New results must be regenerated through the shared target, split, and feature-scaling code in \texttt{quant\_pipeline}.

\subsection{Stage-two trading performance comparison}
The historical DQN table is also retired. Its policy used a wider leveraged action space, normalized the signal using test-period statistics, and omitted transaction costs and slippage. The corrected evaluator constrains positions to the same range for every strategy and charges explicit turnover costs; no new performance claim is made until the policy is retrained under that contract.

The historical stage-two experiment illustrates how a DQN can consume a forecasting signal, but it is not a valid current performance claim under the corrected protocol. The trading environment remains a simplified research simulator and should be interpreted accordingly.

\section{Limitations and Future Work}

There are several limitations in our current project. First, the trading environment is still a simplified return-level simulator. The corrected evaluator includes configurable transaction costs and slippage, but it is not a full order-book or market-impact model. Therefore, any result remains evidence about a research simulation rather than proof that a strategy would work in real trading [4, 5] .

Second, the current model does not consider applying different method in different situations. For example, financial markets behave differently during bull markets, bear markets, and high volatility periods. A single forecasting model and trading agent may not perform well under complicated market conditions. In future work, the model can be improved by training separate strategies for different market conditions.

Third, the DQN's action space is discrete. Although this implementation makes the reinforcement learning problem easier to train and test, the real world portfolio allocation is usually continuous. Future work can try to develop a continuous action reinforcement learning methods, such as policy gradient.

Finally, the current project uses a two stage model structure, where the forecasting model is trained first and then the trading agent uses its prediction signal later. A possible future direction is to develop an end to end reinforcement learning framework that directly optimizes trading performance while learning useful market representations at the same time. 

\section{Conclusion}

In conclusion, this project developed a two stage machine learning framework for stock market direction prediction and trading decision making. The historical results motivate the architecture, but are not current apples-to-apples evidence. The repository now provides a leakage-aware target, chronological split, train-fitted transforms, and cost-aware evaluation contract; retraining all models under that contract is the next step before drawing performance conclusions.


\section*{References}

\par [1] G. E. P. Box, G. M. Jenkins, G. C. Reinsel, and G. M. Ljung, \textit{Time Series Analysis: Forecasting and Control}, 5th ed. Hoboken, NJ, USA: Wiley, 2015.

\par [2] S. Hochreiter and J. Schmidhuber, ``Long short-term memory,'' \textit{Neural Computation}, vol. 9, no. 8, pp. 1735--1780, Nov. 1997. doi:10.1162/neco.1997.9.8.1735

\par [3] V. Mnih, K. Kavukcuoglu, D. Silver, A. A. Rusu, J. Veness, M. G. Bellemare, A. Graves, M. Riedmiller, A. K. Fidjeland, G. Ostrovski, \textit{et al.}, ``Human-level control through deep reinforcement learning,'' \textit{Nature}, vol. 518, no. 7540, pp. 529--533, Feb. 2015. doi:10.1038/nature14236

\par [4] E. F. Fama, ``Efficient capital markets: A review of theory and empirical work,'' \textit{The Journal of Finance}, vol. 25, no. 2, pp. 383--417, May 1970. doi:10.1111/j.1540-6261.1970.tb00518.x

\par [5] J. Moody and M. Saffell, ``Learning to trade via direct reinforcement,'' \textit{IEEE Transactions on Neural Networks}, vol. 12, no. 4, pp. 875--889, Jul. 2001. doi:10.1109/72.935097

\par [6] T. Fischer and C. Krauss, ``Deep learning with long short-term memory networks for financial market predictions,'' \textit{European Journal of Operational Research}, vol. 270, no. 2, pp. 654--669, Oct. 2018. doi:10.1016/j.ejor.2017.11.054

\end{document}
